\section{Empirical Predictions and Regional Policy Interpretation}
\label{sec:policy}

The model is intended as a mechanism-based regional science framework rather than as a literal description of any single hub program. Its empirical value lies in clarifying which observable margins correspond to the theoretical objects introduced above, and in identifying predictions that distinguish \emph{hub-alone coordination}, \emph{pipeline-alone renewal}, and the \emph{bundle} that combines the two. The relevant empirical baseline is therefore richer than a simple hub-versus-no-hub comparison.

A convenient starting point is the decomposition from Sections~\ref{sec:short_run} and~\ref{sec:refresh}. The decentralized support criterion for adding a hub is
\[
\Delta_P^*=S_A-F+\bigl(R_1^P-R_0^P\bigr),
\]
while the planner's criterion is
\[
\Delta_W^*=\Delta_P^*+D_A+M_A.
\]
Here, $S_A$ is the immediate local coordination gain, $D_A$ is the dynamic local correction, and $M_A$ is the additional social value of renewal that decentralized actors do not fully internalize. Likewise, Section~\ref{sec:refresh} shows that
\[
W^*(0)-W(0,0)=R_0>0
\]
for pipeline policy alone, while
\[
W^*(1)-W^*(0)=\Delta_H+\bigl(R_1-R_0\bigr)
\]
determines whether the full bundle is preferred to pipeline alone. These decompositions immediately suggest what should be measured empirically.

\subsection{Mapping Theory to Observable Margins}

Table \ref{tab:proxy_mapping} summarizes a parsimonious mapping from the model's theoretical objects to observable empirical margins.

\begin{table}[htbp]
\centering
\caption{Illustrative Mapping from Theoretical Objects to Empirical Proxies}
\label{tab:proxy_mapping}
\begin{tabular}{p{0.25\textwidth} p{0.27\textwidth} p{0.38\textwidth}}
\toprule
Theoretical object & Empirical interpretation & Illustrative proxies \\
\midrule
Local coordination gain $m_A^1-m_A^0$
& Immediate effect of the hub on within-region matching
& Interaction volume, first-time meetings, new project starts, event-to-project conversion \\

Dynamic redundancy $\delta$
& Speed at which repeated local interaction exhausts novelty
& Repeated collaborator ratio, new ties share, decline in participant turnover, concentration of repeated ties \\

Renewal channel $(\xi,\eta_0,\eta_1,b)$
& Extent to which outside brokerage renews the regional opportunity set and creates direct outside links
& Outside-region collaborator share, invited outside participants, interregional project share, outside mentors, visiting researchers \\

Incremental bundle margin $R_1-R_0$
& Extra value of adding a hub to a renewal policy
& Difference in medium-run novelty, outside collaboration, and repeated-tie dynamics between bundled and pipeline-only settings \\

Absorptive capacity / anchor presence
& Capacity to convert contact and renewal into productive collaboration
& Anchor participation share, university presence, lead-firm involvement, boundary-spanning intermediaries \\
\bottomrule
\end{tabular}
\end{table}

The purpose of Table \ref{tab:proxy_mapping} is not to impose a single empirical design. Rather, it highlights that the model's core variables are observable at least in approximate form. In particular, a credible evaluation should include not only gross activity measures, but also indicators of novelty, repetition, and outside connectivity.

Some proxies are more tightly linked to particular parameters than others. First-time meetings, event-to-project conversion, and the share of new ties speak most directly to the hub-induced matching increment $m_A^1-m_A^0$. Repeated collaborator ratio and repeated-tie concentration are the closest empirical analogues of the redundancy parameter $\delta$. Outside-link conversion and the persistence of new local combinations after outside contact are the most relevant proxies for the renewal channel, with the former most closely related to $\eta_1$ and the latter to $\xi$.

\subsection{An Institutional Anchor}

One reason to take the bundle logic seriously is that current policy practice already contains programs that combine a local support node with organized cross-regional access. A useful illustration is the European Digital Innovation Hubs (EDIH) network. The European Commission describes EDIHs as one-stop shops with a strong regional presence, but also emphasizes that each hub belongs to a wider pan-European network through which firms can access specialized expertise and services located outside their home region \citep{EuropeanCommission2025EDIH}. The institutional architecture is therefore explicitly two-layered: a regionally embedded access point plus network-based referral beyond the home region. In institutional terms, this is close to the distinction developed in the model: a local node reduces coordination frictions inside a region, while the wider network expands the region's effective outside pool.

Of course, EDIHs are not literal copies of the present model. They are digitally oriented service infrastructures, not generic co-creation hubs, and their objectives are broader than the stylized welfare problem analyzed here. Nor do they by themselves validate the model's welfare ranking. But as an institutional fact they are useful because they show that policy practice already treats \emph{local coordination} and \emph{cross-regional access} as complementary rather than interchangeable. That is exactly the policy architecture the model is designed to clarify.

\subsection{Three Testable Predictions}

\noindent\textbf{Prediction 1. Pipeline-only regions should show more outside connectivity than no-policy regions, but bundled regions should show more durable novelty than pipeline-only regions.}

Section~\ref{sec:refresh} shows that pipeline policy alone has positive value, summarized by $R_0>0$. This implies that even without a hub, a region using active brokerage should display more outside-region collaboration than an otherwise similar no-policy region.

However, the bundle is preferred to pipeline alone only when
\[
\Delta_H+\bigl(R_1-R_0\bigr)\ge 0.
\]
This generates the paper's most identifying empirical comparison. Relative to pipeline-only regions, bundled regions should exhibit not just more outside collaboration, but also a more durable novelty profile: lower growth in repeated-tie concentration, higher persistence of new combinations, and stronger conversion of outside links into productive local collaboration.

\noindent\textbf{Prediction 2. Hub creation raises visible local activity, but it need not preserve novelty over time.}

The model implies an immediate coordination gain
\[
B_1(1)-B_1(0)=v_LP_A\bigl(m_A^1-m_A^0\bigr)>0.
\]
Thus, after hub establishment, one should expect more meetings, more matches, and more observed local project activity. However, the dynamic correction
\[
D_A
=
\beta v_LP_A\bigl(m_A^1-m_A^0\bigr)\Bigl[1-\delta\bigl(m_A^1+m_A^0\bigr)\Bigr]
\]
need not be positive. When local redundancy is strong, early activity gains can coexist with a later rise in repeated collaborator ratios and a decline in the share of genuinely new ties.

The empirical implication is straightforward: evaluations based only on attendance, occupancy, or total interactions capture the visible coordination term but not whether novelty persists.

\noindent\textbf{Prediction 3. The bundle advantage should be larger where outside renewal is easier and where absorptive capacity is higher.}

The term $R_1-R_0$ is increasing when outside renewal becomes more valuable. In the baseline model, this occurs when the outside pool is larger, brokerage is more effective, or renewal contributes more strongly to future local novelty. In the extension, the same logic is strengthened by higher absorptive capacity.

The empirical implication is that the same hub policy should look more favorable in regions with richer outside partner pools, stronger interregional brokerage institutions, or stronger anchor actors. Conversely, in places where outside renewal is weak, pipeline policy may still be worthwhile, but the incremental case for adding a hub should be weaker.

These three predictions are not equally diagnostic. Prediction 1 is the paper's most identifying implication because it distinguishes \emph{pipeline alone} from the \emph{bundle} and therefore matches the model's central policy comparison. Prediction 2 is easier to confirm empirically, but it is also less discriminating because many hub programs generate visible local activity. Prediction 3 is best interpreted as a moderator prediction: it identifies the regional conditions under which the bundled strategy should outperform pipeline policy alone by more.

\subsection{Illustrative Parameter Regions}

Figure \ref{fig:rrr_policy_regions} visualizes the model's main comparative policy implication. The dashed line is the hub-alone threshold $F_H^*$ from Section~\ref{sec:hub_efficiency}. The solid curve is the bundled threshold
\[
F_{H+b}^{\dagger}(N_B)=F_H^*(N_A)+\bigl(R_1(N_B)-R_0(N_B)\bigr)
\]
from Section~\ref{sec:refresh}. For fixed costs between these two loci, a hub alone is not justified, but a hub embedded in a renewal strategy is preferred to pipeline policy alone.

\begin{figure}[htbp]
\centering
\begin{tikzpicture}
\begin{axis}[
    width=0.84\textwidth,
    height=0.54\textwidth,
    xmin=0, xmax=4.0,
    ymin=0, ymax=18,
    axis lines=left,
    xlabel={Outside-pool size $N_B$},
    ylabel={Hub fixed cost $F$},
    xtick={0,1,2,3,4},
    ytick={0,3,6,9,12,15,18},
    clip=false,
    legend style={draw=none, fill=none, font=\small, at={(0.02,0.98)}, anchor=north west},
    every axis plot/.append style={thick},
]

\def\PA{15}
\def\vL{1}
\def\mzero{0.15}
\def\mone{0.55}
\def\betaPar{0.9}
\def\deltaPar{0.70}
\def\xiPar{0.10}
\def\etaZero{0.05}
\def\etaOne{0.10}
\def\vX{1}
\def\NA{6}
\def\kappaPar{1}

\pgfmathsetmacro{\DeltaM}{\mone-\mzero}
\pgfmathsetmacro{\SigmaM}{\mone+\mzero}
\pgfmathsetmacro{\FH}{\vL*\PA*\DeltaM*((1+\betaPar)-\betaPar*\deltaPar*\SigmaM)}

\addplot[name path=Fbundle, domain=0:4, samples=200, black]
    {\FH + ((\betaPar^2)/(2*\kappaPar))*(((\xiPar*\vL*\PA*\mone)+(\vX*\NA*x*(\etaZero+\etaOne)))^2 - ((\xiPar*\vL*\PA*\mzero)+(\vX*\NA*x*\etaZero))^2)};
\addlegendentry{$F_{H+b}^{\dagger}(N_B)$}

\addplot[name path=Fhub, domain=0:4, samples=2, dashed, black]
    {\FH};
\addlegendentry{$F_H^*$ (hub alone)}

\addplot[name path=zero, domain=0:4, samples=2, draw=none] {0};
\addplot[gray!18, draw=none] fill between[of=Fhub and zero];
\addplot[gray!45, draw=none] fill between[of=Fbundle and Fhub];

\node[font=\small, align=center] at (axis cs:1.0,4.4)
    {Hub creation justified\\even without renewal};

\node[font=\small, align=center] at (axis cs:2.35,10.8)
    {Hub alone not justified,\\but bundle preferred\\to pipeline alone};

\node[font=\small, align=center] at (axis cs:2.9,15.6)
    {Pipeline alone preferred};

\end{axis}
\end{tikzpicture}
\caption{Illustrative policy regions in $(N_B,F)$ space. The dashed line is the hub-alone threshold $F_H^*$. The solid curve is the bundled threshold $F_{H+b}^{\dagger}(N_B)=F_H^*+\bigl(R_1(N_B)-R_0(N_B)\bigr)$. Between the two loci, hub creation is justified only when embedded in a renewal strategy. Parameter values are illustrative: $N_A=6$, $v_L=v_X=1$, $m_A^0=0.15$, $m_A^1=0.55$, $\beta=0.9$, $\delta=0.70$, $\xi=0.10$, $\eta_0=0.05$, $\eta_1=0.10$, and $\kappa=1$.}
\label{fig:rrr_policy_regions}
\end{figure}

Figure \ref{fig:rrr_policy_regions} should be read only as an illustration of the maintained functional form. It is not a calibration, an estimate, or a policy frontier recovered from data. Its role is simply to visualize how the conditional inequality in \eqref{eq:bundle_vs_pipeline} partitions the $(N_B,F)$ space under the benchmark assumptions.

\subsection{Implications for Empirical Design}

Taken together, the theoretical decomposition, the proxy mapping, and Figure \ref{fig:rrr_policy_regions} imply that a theory-consistent empirical evaluation of hub policy should be dynamic, network-aware, and explicitly regional.

First, evaluation should be \emph{dynamic}. Early activity measures capture the immediate coordination gain, but do not show whether novelty persists. At a minimum, researchers should compare initial interaction gains with the subsequent evolution of new-ties share, repeated collaborator ratio, and participant turnover.

Second, evaluation should be \emph{network-aware}. The model is fundamentally about the difference between creating interactions and creating \emph{nonrepetitive} interactions. Measures such as repeated collaborator ratio, outside-region collaborator share, and the novelty content of joint applications or joint projects are therefore better aligned with the theory than aggregate attendance alone.

Third, evaluation should be \emph{explicitly regional}. The outside pool is not a residual background feature. It is a policy-relevant object. Data on external participants, outside mentors, nonlocal co-applicants, or interregional partner flows are central outcome variables for judging whether a hub is functioning as part of a broader renewal strategy.

A minimal empirical design therefore requires more than a treated-versus-untreated comparison. The most informative setting is one in which researchers can distinguish among three baselines: no organized intervention, pipeline-oriented renewal without a hub, and the bundled strategy that combines local coordination with outside brokerage. In such a design, Prediction 1 is the most discriminating test. If the model is right, the strongest evidence for the bundled logic is not merely that hub regions become busier, but that bundled regions outperform pipeline-only regions on outside-link conversion, repeated-tie dynamics, and the persistence of genuinely new combinations.

In practice, the cleanest empirical settings would exploit staggered introductions of hub programs and interregional brokerage schemes, or program-level administrative data that distinguish local-only support from collaborative networking support. Panel event-study or matched-cohort designs could then trace not only activity effects, but also changes in outside-link conversion, repeated-tie concentration, and the persistence of new combinations across the three policy baselines. A minimal reduced-form specification would interact a bundle indicator with post-treatment time and test whether the coefficient on the bundle, relative to pipeline-only support, is positive for outside-link conversion and negative for repeated-tie concentration. This emphasis is consistent with the collaborative-policy literature: policy-induced R\&D networks do not automatically generate stronger inter-regional knowledge diffusion \citep{BednarzBroekel2019}, and collaborative subsidies appear most relevant when they shift regions toward new technological combinations rather than merely increasing subsidized activity \citep{MewesBroekel2020}.

Overall, the model points toward a more demanding empirical agenda than a conventional utilization-based evaluation. The critical question is not simply whether the hub generated more activity, but whether it generated activity that remained novel, whether pipeline policy created outside renewal even without a hub, and whether the bundle outperformed that pipeline-alone baseline.
